\documentclass[11pt, oneside]{article} 
\usepackage{geometry}
\geometry{letterpaper} 
\usepackage{graphicx}
	
\usepackage{amssymb}
\usepackage{amsmath}
\usepackage{parskip}
\usepackage{color}
\usepackage{hyperref}

\graphicspath{{/Users/telliott/Dropbox/Github-Math/figures/}}
% \begin{center} \includegraphics [scale=0.4] {gauss3.png} \end{center}

\title{Proof}
\date{}

\begin{document}
\maketitle
\Large

%[my-super-duper-separator]

Here's a challenging problem from Coxeter that can be solved by an ingenious construction.

Let $ABCD$ be a parallelogram and let $P$ be found such that the angles labeled $\alpha$ are equal:  $\angle DPA = \angle BPC$.

\begin{center} \includegraphics [scale=0.35] {Coxeter_hard1.png} \end{center}

To prove:  the angles labeled $\theta$ and $\phi$ are also equal:  $\angle DPC = \angle BPC$.

\emph{Proof}.

Like many problems, the solution starts with an inspired construction.  

Find $Q$ such that $PQ \parallel CB$. 

\begin{center} \includegraphics [scale=0.35] {Coxeter_hard2.png} \end{center}

As a result, we have two new parallelograms $BQPC$ and $AQPD$.

It also gives us new angles equal to $\alpha$, $\theta$ and $\phi$

\begin{center} \includegraphics [scale=0.35] {Coxeter_hard2a.png} \end{center}

And that may suggest, to the prepared mind, that we need to find a circle with all four points $A,B,Q,P$ lying on it.

Draw the circumcircle for $\triangle ABQ$.

\begin{center} \includegraphics [scale=0.35] {Coxeter_hard4b.png} \end{center}

$BQPC$ is a parallelogram.  The diagonal forms congruent triangles so $\angle BPQ$ is also equal to $\alpha$ and it is subtended by arc $BQ$ on the circle.

$\angle BAQ$ is formed by two lines parallel to the lines forming $\angle PDC$, so it also has measure $\alpha$ and is subtended by $BQ$.

By the converse of the cyclic quadrilateral theorem, it follows that $P$ is on the circle.

Using diagonals of the parallelograms, we find $\angle PAQ = \phi$ and $\angle PBQ = \theta$.

But both angles are subtended by the same arc of the circle $PQ$, so they are equal.

$\square$.

\end{document}